\section{一个地址怎样在 RAM 中选中确切的字节}

前文把 RAM 画成一个方框，并规定它占用一段物理地址。这个方框还需要展开到足以解释
“为什么地址加一会得到下一个字节”。本节不重复存储电路的门级实现，只说明控制一次访问
所必需的内部结构。

\subsection{容量来自地址数量与每个位置的宽度}

本章 RAM 范围为 \code{0x00100000--0x001fffff}，共
\[
\texttt{0x001fffff}-\texttt{0x00100000}+1=\texttt{0x00100000}
\]
个字节，即 \(2^{20}\) 字节。RAM 内部不需要保存完整的高位地址。互连已经证明请求落在
这个范围后，向 RAM 提供范围内偏移：
\[
\texttt{offset}=\texttt{physical\_address}-\texttt{0x00100000}.
\]
偏移的低 20 位足以选择 \(2^{20}\) 个字节位置。

\begin{center}
\begin{tikzpicture}[node distance=11mm and 15mm]
\node[stage,minimum width=3.1cm] (addr) {物理地址\\32 位};
\node[stage,minimum width=3.2cm,right=of addr] (range) {互连范围比较\\高位决定目标};
\node[stage,minimum width=3.2cm,right=of range] (offset) {减去 RAM 基址\\得到 20 位偏移};
\node[stage,minimum width=3.2cm,below=14mm of offset] (array) {RAM 字节阵列\\选中一个或四个相邻字节};
\draw[flow] (addr) -- (range);
\draw[flow] (range) -- (offset);
\draw[flow] (offset) -- (array);
\end{tikzpicture}
\end{center}
\diagramnote{完整物理地址用于整机选择目标；目标内部只需使用足够区分自身位置的偏移位。}

如果直接把 \(2^{20}\) 个字节画成一条线，硬件结构会显得不现实。实际阵列通常分成行、列
和若干存储体。对本章的软件契约而言，这些组织方式可以不同，只要同一地址总返回同一组
已提交字节，并遵守规定的响应次序。

\subsection{为何一次四字节读取仍以字节地址命名}

处理器的地址单位是字节。\code{LOAD} 宽度为四时，地址 \(A\) 表示读取
\(A,A+1,A+2,A+3\) 四个相邻字节。地址没有自动除以四；只是 RAM 可以把低两位用来选择
一个四字节组中的起点。

\begin{longtable}{|L{0.20\linewidth}|L{0.25\linewidth}|L{0.43\linewidth}|}
\hline
\rowcolor{BookTableHead}
地址低两位 & 四字节访问是否对齐 & 本章处理方式 \\ \hline
\code{00} & 是 & 读取或写入同一个四字节组 \\ \hline
\code{01} & 否 & 返回 \code{ALIGNMENT}，不改变任何字节 \\ \hline
\code{10} & 否 & 返回 \code{ALIGNMENT}，不改变任何字节 \\ \hline
\code{11} & 否 & 返回 \code{ALIGNMENT}，不改变任何字节 \\ \hline
\end{longtable}

禁止未对齐访问不是数学必然。某些处理器会拆成两次阵列读取，再在内部拼接。本章选择拒绝，
是为了让每条四字节事务只涉及一个对齐组，错误时也容易保证“全写或全不写”。程序因此
必须让整数地址和 SP 都是四的倍数。

\subsection{四条字节通道怎样组成一个 32 位值}

可以把数据端口理解为四条八位通道。对齐地址 \(A\) 的低两位为零；第零通道连接 \(A\)，
第一通道连接 \(A+1\)，依此类推。读取时，小端组合器产生：
\[
\texttt{value}=
b_0+(b_1\ll8)+(b_2\ll16)+(b_3\ll24).
\]

\begin{center}
\begin{tikzpicture}[x=1.55cm,y=0.82cm]
\foreach \x/\lab in {0/{\(A\)},1/{\(A+1\)},2/{\(A+2\)},3/{\(A+3\)}} {
  \draw[fill=BookTint,draw=BookRule] (\x*1.7,2.2) rectangle ++(1.3,0.8);
  \node[font=\footnotesize] at (\x*1.7+0.65,2.6) {\lab};
  \draw[->,>=Latex,thick,BookAccent] (\x*1.7+0.65,2.2) -- (\x*1.7+0.65,1.35);
}
\draw[rounded corners=2pt,fill=BookAccent!12,draw=BookRule] (0.15,0.65) rectangle (6.25,1.35);
\node[font=\footnotesize] at (3.2,1.0) {小端组合器：按地址次序把四个字节放入四条八位通道};
\draw[->,>=Latex,thick,BookAccent] (3.2,0.65) -- (3.2,-0.05);
\draw[rounded corners=2pt,fill=BookTint,draw=BookRule] (1.7,-0.75) rectangle (4.7,-0.05);
\node[font=\footnotesize] at (3.2,-0.4) {32 位返回值 \(b_0+(b_1\ll8)+(b_2\ll16)+(b_3\ll24)\)};
\end{tikzpicture}
\end{center}
\diagramnote{字节顺序是处理器与存储接口之间的解释规则。RAM 单元只保存八个位，并不保存
“最低有效字节”这个语义标签。}

\subsection{写入时为什么需要字节使能}

写请求除了 32 位数据，还带四位 \code{byte\_enable}。四字节 \code{STORE} 使用
\code{1111}；若以后增加单字节存储，地址低位和字节使能可只选择一条通道。例如在地址
\(A+2\) 写入 \code{0x7f}，可能使用：

\begin{lstlisting}
aligned_address = A
write_data      = 0x007f0000
byte_enable     = 0100
\end{lstlisting}

RAM 必须先同时检查范围、对齐规则和字节使能是否合法，再更新被选择的通道。若先更新第零
通道、随后才发现第四个字节越界，软件将面对部分写入。本章接口把检查与更新分开，成功边沿
一次提交全部被使能字节。

\subsection{读取端口和写入端口是否能同时工作}

当前机器只有一个处理器请求发起者，互连一次只接受一笔未完成事务。RAM 因而使用一个共享
端口就足够：某个周期接受读或写，经过固定或可变延迟后给出一次响应。在响应出现以前，
请求字段必须由发起者或 RAM 内部锁存。

\begin{lstlisting}
ram_pending:
  valid
  operation
  offset
  width
  write_data
  byte_enable
\end{lstlisting}

这些字段解释了“RAM 正在忙”究竟保存什么。若只保存地址而忘记操作方向，响应时就无法
判断应返回数据还是提交写入；若不保存字节使能，等待期间输入线变化会改变原请求含义。

\begin{longtable}{|L{0.18\linewidth}|L{0.29\linewidth}|L{0.41\linewidth}|}
\hline
\rowcolor{BookTableHead}
内部阶段 & 已锁存的状态 & 对外接口状态 \\ \hline
\code{IDLE} & 无待处理请求 & \code{req\_ready=1} \\ \hline
\code{CHECK} & 请求全部字段 & 暂不接受第二笔请求 \\ \hline
\code{READ} & 偏移与宽度 & 阵列选择相应字节 \\ \hline
\code{WRITE} & 偏移、数据、字节使能 & 成功边沿提交所选字节 \\ \hline
\code{RESPOND} & 数据或错误码 & \code{resp\_valid=1} \\ \hline
\end{longtable}

阶段名只是本章的一种实现。软件不能读取这些名称，只能观察到请求是否被接受、响应何时
有效以及成功响应对应什么效果。把内部状态列出，是为了证明接口承诺有足够的物理状态可以
实现，而不是把特定状态机强加给所有 RAM。

\subsection{启动存储与 RAM 为什么不能合成同一种可写阵列}

如果复位入口也位于普通可写 RAM，操作者准备任务时的一次错址就可能覆盖下一次启动所需的
第一条指令。机器随后连检查错误记录的代码都无法取得。将启动程序放在普通运行不能写的
存储中，保留了一个比任务内容更稳定的起点。

两种存储仍可共享读取请求格式：

\begin{longtable}{|L{0.20\linewidth}|L{0.28\linewidth}|L{0.40\linewidth}|}
\hline
\rowcolor{BookTableHead}
性质 & 启动存储 & RAM \\ \hline
复位后内容 & 制造或维护阶段已确定 & 本章不规定 \\ \hline
普通读取 & 支持 & 支持 \\ \hline
普通写入 & 返回 \code{READ\_ONLY} & 合法范围内提交 \\ \hline
本章用途 & 检查、交接、返回入口 & 任务代码、数据、栈与工作区 \\ \hline
\end{longtable}

因此“都能返回字节”不等于“具有相同生命周期”。启动存储给出稳定起点，RAM 给出运行期
可变状态；地址互连把同一种请求送到承担不同责任的目标。

\subsection{物理地址表仍然不是内存保护表}

范围比较只回答地址属于哪个器件。任务向 \code{0x001ff000} 写入时，互连会正确选择 RAM，
RAM 也会正确更新准备记录；两个器件都履行了当前契约，却破坏了下一次启动依据。要拒绝这
笔写入，判断还需要知道“当前请求者身份”和“该范围允许何种访问”。第一章尚未建立这两项
输入，所以不能从地址表推断保护已经存在。

\subsection{复位为什么不顺便清空全部 RAM}

处理器复位只需设置少量控制寄存器，几个边沿内即可完成。RAM 有 \(2^{20}\) 个字节；
若要求复位信号直接把每个存储单元同时改成零，就要为整个阵列增加清零通路，并规定清零
期间普通访问如何处理。另一种实现是由启动程序逐地址写零，但这需要发出
\(2^{18}\) 笔四字节写，复位释放到首次检查之间会出现很长的执行阶段。

本章没有任何需求要求全部 RAM 为零。代码区、输入区、结果槽、返回栈字和准备记录都会由
操作者明确写入；启动程序工作区只在首次写入后读取。其余字节保持未规定更简单，也迫使
程序不依赖自己尚未建立的初值。

\begin{longtable}{|L{0.22\linewidth}|L{0.29\linewidth}|L{0.37\linewidth}|}
\hline
\rowcolor{BookTableHead}
可选策略 & 获得的性质 & 付出的代价或仍有的限制 \\ \hline
硬件全阵列清零 & 复位后每个 RAM 字节为零 & 阵列增加清零结构，复位时序更复杂 \\ \hline
启动程序逐字清零 & 无需阵列专用清零线 & 启动时间随 RAM 容量增长 \\ \hline
只初始化将要使用的区域 & 启动工作与本次需求成比例 & 未声明区域必须始终视为未规定 \\ \hline
\end{longtable}

第三种策略形成一条重要纪律：任何代码在读取某个可写字节前，都要能指出本次运行中哪一步
先写入了它。若回答只能追溯到“机器通常会给零”，该读取就没有建立在本章契约上。

\subsection{调试夹具怎样避免端口争用}

人工准备要求外部夹具能够访问 RAM，而正常运行时 RAM 端口属于处理器。两者若同时驱动
地址线，电气上会冲突，协议上也无法判断哪笔请求先发生。本章利用复位状态设置一个简单
选择器：

夹具内部至少需要保存当前地址、访问方向、一个字节的写数据以及尚未完成标志。操作者给出
一项读写命令后，命令接收级把这些字段写入请求寄存器；事务控制器保持字段不变，直到 RAM
返回成功或错误；读操作得到的字节写入读回寄存器，随后才向操作者报告完成。若夹具只把
按钮或主机信号直接接到 RAM 引脚，输入在事务期间改变时便无法证明实际访问了哪个地址。

\begin{center}
\begin{tikzpicture}[node distance=9mm and 6mm]
\node[stage,minimum width=2.65cm] (host) {操作者或主机\\字节、地址与命令};
\node[stage,minimum width=2.75cm,right=of host] (command)
  {命令接收级\\检查字段是否完整};
\node[stage,minimum width=2.85cm,right=of command] (registers)
  {请求寄存器\\地址、方向、写数据};
\node[stage,minimum width=2.75cm,right=of registers] (control)
  {事务控制器\\保持请求直到响应};
\node[stage,minimum width=2.7cm,below=12mm of control] (adapter)
  {RAM 端口适配\\请求与响应信号};
\node[stage,minimum width=2.85cm,left=of adapter] (readback)
  {读回与状态寄存器\\数据、错误、完成};
\draw[flow] (host) -- (command);
\draw[flow] (command) -- (registers);
\draw[flow] (registers) -- (control);
\draw[flow] (control) -- (adapter);
\draw[flow] (adapter) -- (readback);
\draw[flow] (readback.west) -| (host.south);
\end{tikzpicture}
\end{center}
\diagramnote{夹具的提交边界是 RAM 响应到达并写入状态寄存器。主机发出命令并不等于 RAM
已经改变；读回核对也必须在对应事务完成以后进行。}

\begin{center}
\begin{tikzpicture}[node distance=11mm and 16mm]
\node[stage,minimum width=3.0cm] (cpu) {处理器请求\\仅在复位释放后有效};
\node[stage,minimum width=3.0cm,below=14mm of cpu] (fixture) {调试夹具请求\\仅在复位保持时有效};
\node[stage,minimum width=3.2cm,right=22mm of cpu,yshift=-7mm] (mux) {端口选择器\\由复位状态控制};
\node[stage,minimum width=3.0cm,right=of mux] (ram) {RAM 共享端口};
\draw[flow] (cpu) -- (mux);
\draw[flow] (fixture) -- (mux);
\draw[flow] (mux) -- (ram);
\end{tikzpicture}
\end{center}
\diagramnote{选择器不是运行中的第二个请求发起者。复位状态保证同一时刻只有一侧能够产生
有效请求，因此第一章仍是单发起者机器。}

夹具写入期间，处理器被复位阻止，不能观察半成品；复位释放以后，选择器固定选择处理器，
夹具即使改变自己的地址输入也不会触及 RAM。操作者若要再次修改内容，必须重新保持复位，
先让处理器停止普通事务，再取得端口。

\begin{longtable}{|L{0.20\linewidth}|L{0.22\linewidth}|L{0.22\linewidth}|L{0.24\linewidth}|}
\hline
\rowcolor{BookTableHead}
复位状态 & 处理器端口 & 夹具端口 & RAM 接受者 \\ \hline
保持有效 & 不发普通请求 & 可读写 & 夹具 \\ \hline
正在释放的规定边沿 & 仍不接受新请求 & 结束最后一笔请求 & 无；完成所有权转换 \\ \hline
已经释放 & 可发普通请求 & 被隔离 & 处理器 \\ \hline
\end{longtable}

中间一行防止夹具最后一笔写尚未响应，处理器就开始第一笔取指。复位控制器只有在夹具报告
空闲后才允许释放；因此 \code{prepared=1} 的写响应先于第一条启动指令。这个硬件次序与
前文的人工提交顺序共同保证，启动程序不会读到正在变化的准备记录。
